\section{Discussion}
\label{sec:discussion}

\subsection{Atmospheric Characterization Potential}
\label{sec:atmos}

With $\mathrm{TSM} = 94 \pm 12$, TOI-4521\,b satisfies the
\citet{Kempton2018} threshold of TSM $> 92$ for high-priority
atmospheric follow-up in the sub-Neptune regime.  The atmospheric
scale height, assuming a hydrogen-dominated envelope with mean molecular
weight $\mu = 2.3\,\mathrm{amu}$ and adopting $g_p$ from a mass-radius
relation, is $H \approx 200\,\mathrm{km}$, corresponding to a
single-scale-height spectral feature amplitude of
$\Delta\delta \simeq 2H R_p / R_\star^2 \approx 80\,\mathrm{ppm}$.
This is comfortably detectable with $\sim 20\,\mathrm{hr}$ of NIRSpec/G395H
integration.

By contrast, a high-mean-molecular-weight atmosphere ($\mu = 28\,\mathrm{amu}$,
analogous to Venus) would produce features of only $\sim 6\,\mathrm{ppm}$,
effectively featureless at \textit{JWST} precision.  TOI-4521\,b therefore
provides a decisive discriminator: a detection of spectral features
would strongly favor a hydrogen-dominated (and potentially Hycean)
composition, while a flat spectrum would indicate atmospheric loss or
a high-$\mu$ envelope.

\subsection{Companion Star Contamination}
\label{sec:companion}

The neighboring K4\,V star at $\rho = 1\farcs3$ contributes a dilution
fraction $f_\mathrm{cont} = F_\mathrm{comp}/(F_\star + F_\mathrm{comp})
= 0.041 \pm 0.008$ in the \textit{TESS} bandpass.  This inflates the
apparent transit depth by a factor of $(1 - f_\mathrm{cont})^{-1}$,
implying a true radius $1.4\%$ larger than the uncorrected measurement.
All reported parameters in Table~\ref{tab:params} include this correction.
High-contrast imaging with Gemini-North/'Alopeke confirmed the companion
is not physically associated (proper-motion baseline of 18\,yr from
\textit{Gaia} and 2MASS astrometry), and is therefore a foreground
interloper.

%% -------------------------------------------------------
%%  8. Conclusions
%% -------------------------------------------------------
